% GENERATED FILE. Edit canonical lessons, then run npm run generate:books.
% canonical-source-hash: fnv1a64:1e597219207a5ef8
% canonical-lessons: UR-C348-etiraz, UR-C348-tanqid, UR-C348-nuqta-e-nazar, UR-C348-dhamkana, UR-C348-phuslana

\chapter{Objecting and Persuading}
\label{ch:ur-objecting-and-persuading}
\hlchaptermodality{348}

\begin{chapteropening}
\textbf{By the end of this chapter:} \emph{I can object, criticise, threaten or coax.}

Complete the last lesson of chapter 348: I can object, criticise, threaten or coax.
\end{chapteropening}

\section[\texorpdfstring{\ur{اعتراض} (etirāz)}{اعتراض (etirāz)}]{\ur{اعتراض} (etirāz) — an objection}

\label{lesson:UR-C348-etiraz}



\begin{quote}
Before the new one: say the Urdu for to get tangled, then the Urdu for to bring up a topic.
\end{quote}

\subsection*{You'll want to know: \ur{اعتراض}}
\textbf{\ur{اعتراض}} — \emph{etirāz} — \textquotedblleft{}an objection\textquotedblright{}.

\begin{quote}
\ur{کیا} \ur{آپ} \ur{کو} \ur{کوئی} \ur{اعتراض} \ur{ہے؟}
\end{quote}

\begin{quote}
\emph{kyā āp ko koī etirāz hai?}
\end{quote}

\begin{quote}
Do you have any objection?
\end{quote}

Masculine.

\begin{cousinweb}
\textbf{\ur{اعتراض}} (\emph{etirāz}) is usually traced, through Persian, to Arabic \emph{iʿtirāḍ}, \textquotedblleft{}standing across the way; objecting\textquotedblright{}.
\end{cousinweb}

\begin{grammarlens}[title={What you've built}]
One more word for giving reasons and opinions.
\end{grammarlens}

\subsection*{Your turn}
\begin{itemize}
\raggedright
  \item \emph{Say it:} \emph{etirāz}
  \item \emph{Say it:} \emph{etirāz}, once more
  \item \emph{Say it:} \emph{kyā āp ko koī etirāz hai?} — Do you have any objection?
  \item \emph{Recall:} say the Urdu for to get tangled, then the Urdu for to bring up a topic, then say \emph{etirāz} again
\end{itemize}

\begin{tcolorbox}[breakable,colback=teal!4,colframe=teal!35!black,title={Before you move on}]
What does \ur{اعتراض} mean? (\textquotedblleft{}An objection\textquotedblright{}.) Say it once more.
\end{tcolorbox}
\section[\texorpdfstring{\ur{تنقید} (tanqīd)}{تنقید (tanqīd)}]{\ur{تنقید} (tanqīd) — criticism}

\label{lesson:UR-C348-tanqid}



\begin{quote}
Before the new one: say the Urdu for to bring up a topic, then the Urdu for an objection.
\end{quote}

\subsection*{You'll want to know: \ur{تنقید}}
\textbf{\ur{تنقید}} — \emph{tanqīd} — \textquotedblleft{}criticism\textquotedblright{}.

\begin{quote}
\ur{اس} \ur{نے} \ur{حکومت} \ur{پر} \ur{سخت} \ur{تنقید} \ur{کی۔}
\end{quote}

\begin{quote}
\emph{us ne hukūmat par sakht tanqīd kī.}
\end{quote}

\begin{quote}
He criticised the government harshly.
\end{quote}

Feminine.

\begin{cousinweb}
\textbf{\ur{تنقید}} (\emph{tanqīd}) is usually traced, through Persian, to Arabic \emph{tanqīd}, \textquotedblleft{}sorting good coins from bad\textquotedblright{}: judging what holds up.
\end{cousinweb}

\begin{grammarlens}[title={What you've built}]
One more word for giving reasons and opinions.
\end{grammarlens}

\subsection*{Your turn}
\begin{itemize}
\raggedright
  \item \emph{Say it:} \emph{tanqīd}
  \item \emph{Say it:} \emph{tanqīd}, once more
  \item \emph{Say it:} \emph{us ne hukūmat par sakht tanqīd kī.} — He criticised the government harshly.
  \item \emph{Recall:} say the Urdu for to bring up a topic, then the Urdu for an objection, then say \emph{tanqīd} again
\end{itemize}

\begin{tcolorbox}[breakable,colback=teal!4,colframe=teal!35!black,title={Before you move on}]
What does \ur{تنقید} mean? (\textquotedblleft{}Criticism\textquotedblright{}.) Say it once more.
\end{tcolorbox}
\section[\texorpdfstring{\ur{نقطہ} \ur{نظر} (nuqta-e-nazar)}{نقطہ نظر (nuqta-e-nazar)}]{\ur{نقطہ} \ur{نظر} (nuqta-e-nazar) — a viewpoint; a point of view}

\label{lesson:UR-C348-nuqta-e-nazar}



\begin{quote}
Before the new one: say the Urdu for an objection, then the Urdu for criticism.
\end{quote}

\subsection*{You'll want to know: \ur{نقطہ} \ur{نظر}}
\textbf{\ur{نقطہ} \ur{نظر}} — \emph{nuqta-e-nazar} — \textquotedblleft{}a viewpoint; a point of view\textquotedblright{}.

\begin{quote}
\ur{آپ} \ur{کا} \ur{نقطہ} \ur{نظر} \ur{کیا} \ur{ہے؟}
\end{quote}

\begin{quote}
\emph{āp kā nuqta-e-nazar kyā hai?}
\end{quote}

\begin{quote}
What is your point of view?
\end{quote}

Masculine.

\begin{cousinweb}
\textbf{\ur{نقطہ} \ur{نظر}} (\emph{nuqta-e-nazar}) is built from two loans, each usually traced, through Persian, to Arabic: Arabic \emph{nuqṭa}, \textquotedblleft{}a dot, a point\textquotedblright{}, and Arabic \emph{naẓar}, \textquotedblleft{}a look\textquotedblright{}. The point you look from.
\end{cousinweb}

\begin{grammarlens}[title={What you've built}]
One more word for giving reasons and opinions.
\end{grammarlens}

\subsection*{Your turn}
\begin{itemize}
\raggedright
  \item \emph{Say it:} \emph{nuqta-e-nazar}
  \item \emph{Say it:} \emph{nuqta-e-nazar}, once more
  \item \emph{Say it:} \emph{āp kā nuqta-e-nazar kyā hai?} — What is your point of view?
  \item \emph{Recall:} say the Urdu for an objection, then the Urdu for criticism, then say \emph{nuqta-e-nazar} again
\end{itemize}

\begin{tcolorbox}[breakable,colback=teal!4,colframe=teal!35!black,title={Before you move on}]
What does \ur{نقطہ} \ur{نظر} mean? (\textquotedblleft{}A viewpoint; a point of view\textquotedblright{}.) Say it once more.
\end{tcolorbox}
\section[\texorpdfstring{\ur{دھمکانا} (dhamkānā)}{دھمکانا (dhamkānā)}]{\ur{دھمکانا} (dhamkānā) — to threaten}

\label{lesson:UR-C348-dhamkana}



\begin{quote}
Before the new one: say the Urdu for criticism, then the Urdu for a viewpoint.
\end{quote}

\subsection*{You'll want to know: \ur{دھمکانا}}
\textbf{\ur{دھمکانا}} — \emph{dhamkānā} — \textquotedblleft{}to threaten\textquotedblright{}.

\begin{quote}
\ur{اس} \ur{نے} \ur{دکاندار} \ur{کو} \ur{دھمکایا۔}
\end{quote}

\begin{quote}
\emph{us ne dukāndār ko dhamkāyā.}
\end{quote}

\begin{quote}
He threatened the shopkeeper.
\end{quote}

Past, with \ur{نے} (he or she …): \ur{اس} \ur{نے} \ur{مجھے} \ur{دھمکایا}, \emph{us ne mujhe dhamkāyā}.

\begin{cousinweb}
\textbf{\ur{دھمکانا}} (\emph{dhamkānā}) is inherited: CDIAL lists the Hindi cousin \emph{dhamkī}, \textquotedblleft{}a threat\textquotedblright{}, under \emph{*dhammakka}, a form it reconstructs (CDIAL 6736).
\end{cousinweb}

\begin{grammarlens}[title={What you've built}]
One more word for giving reasons and opinions.
\end{grammarlens}

\subsection*{Your turn}
\begin{itemize}
\raggedright
  \item \emph{Say it:} \emph{dhamkānā}
  \item \emph{Say it:} \emph{dhamkānā}, once more
  \item \emph{Say it:} \emph{us ne dukāndār ko dhamkāyā.} — He threatened the shopkeeper.
  \item \emph{Recall:} say the Urdu for criticism, then the Urdu for a viewpoint, then say \emph{dhamkānā} again
\end{itemize}

\begin{tcolorbox}[breakable,colback=teal!4,colframe=teal!35!black,title={Before you move on}]
What does \ur{دھمکانا} mean? (\textquotedblleft{}To threaten\textquotedblright{}.) Say it once more.
\end{tcolorbox}
\section[\texorpdfstring{\ur{پھسلانا} (phuslānā)}{پھسلانا (phuslānā)}]{\ur{پھسلانا} (phuslānā) — to coax; to wheedle}

\label{lesson:UR-C348-phuslana}



\begin{quote}
Before the new one: say the Urdu for a viewpoint, then the Urdu for to threaten.
\end{quote}

\subsection*{You'll want to know: \ur{پھسلانا}}
\textbf{\ur{پھسلانا}} — \emph{phuslānā} — \textquotedblleft{}to coax; to wheedle\textquotedblright{}.

\begin{quote}
\ur{اس} \ur{نے} \ur{پھسلا} \ur{کر} \ur{ابو} \ur{سے} \ur{پیسے} \ur{لیے۔}
\end{quote}

\begin{quote}
\emph{us ne phuslā kar abbū se paise liye.}
\end{quote}

\begin{quote}
She coaxed some money out of her father.
\end{quote}

Past, with \ur{نے} (he or she …): \ur{اس} \ur{نے} \ur{بچے} \ur{کو} \ur{پھسلایا}, \emph{us ne bachche ko phuslāyā}.

\begin{cousinweb}
\textbf{\ur{پھسلانا}} (\emph{phuslānā}) is inherited: CDIAL lists the Hindi cousin \emph{phuslānā}, \textquotedblleft{}to cajole\textquotedblright{}, under Sanskrit \emph{spṛśáti}, \textquotedblleft{}touches, strokes\textquotedblright{} (CDIAL 13815). To win someone over as one strokes an animal.
\end{cousinweb}

\begin{grammarlens}[title={What you've built}]
One more word for giving reasons and opinions.
\end{grammarlens}

\subsection*{Your turn}
\begin{itemize}
\raggedright
  \item \emph{Say it:} \emph{phuslānā}
  \item \emph{Say it:} \emph{phuslānā}, once more
  \item \emph{Say it:} \emph{us ne phuslā kar abbū se paise liye.} — She coaxed some money out of her father.
  \item \emph{Recall:} say the Urdu for a viewpoint, then the Urdu for to threaten, then say \emph{phuslānā} again
\end{itemize}

\begin{tcolorbox}[breakable,colback=teal!4,colframe=teal!35!black,title={Before you move on}]
What does \ur{پھسلانا} mean? (\textquotedblleft{}To coax; to wheedle\textquotedblright{}.) Say it once more.
\end{tcolorbox}
